\chapter{由反射生成的群；根系。}

本注所考虑的群，是同几何学、分析学以及李群理论的各种问题相关联地出现的，它们有时以置换群的形式出现，有时以欧几里得几何或双曲几何中位移群的形式出现，而这些不同的观点直到晚近才得到协调。

从历史上看，这一理论的开端要比群概念的引入早得多：它的源头其实在于对几何图形的“正则性”或“对称性”的研究，尤其是正多边形与正多面体的确定（无疑可以追溯到毕达哥拉斯学派），后者构成了欧几里得《几何原本》的压卷之作，也是希腊天才最令人赞叹的创造之一。后来，特别是在中世纪盛期的阿拉伯作者那里，然后在开普勒那里，出现了用两两全等（但不必正则）的多边形对平面或球面作正则“铺砌”的数学理论的雏形，这无疑与古代文明和阿拉伯文明所构想的各种类型的装饰的起源有联系（这些装饰有充分的理由被视为这些文明所发展的数学的一个真实组成部分{[}291{]}）。

大约在 1830 至 1840 年间，晶体学方面的研究（赫塞尔、布拉维、默比乌斯）引导人们考虑一个问题，它恰好就是确定三维欧几里得空间中有限位移群的问题，尽管上述作者当时还没有使用群论的语言；群论的语言直到 1860 年前后才逐渐通行，而正是在对群进行分类的背景之下，若尔当于 1869 年{[}174 b{]}确定了 \(\mathbf{R}^3\) 中保持定向的离散位移子群（更一般地，确定了保持定向的位移群的所有闭子群）。

直到十九世纪的最后几年，这股思想潮流都在朝许多方向发展，其中最显著的如下：

\begin{enumerate}
\def\labelenumi{\arabic{enumi}.}
\tightlist
\item
  与有限群理论中出现得很早的一种倾向相一致，人们寻求用简单类型的生成元和关系来给出有限位移群的“表出”。正是这样，哈密顿早在 1856 年{[}145 b{]}就证明了欧几里得空间中
\end{enumerate}

\(R^{3}\) 的有限旋转群由两个生成元 S、T 生成，它们满足关系 \(S^{p} = T^{q} = (ST)^{3} = 1\)，其中 p 与 q 取适当的值。

\begin{enumerate}
\def\labelenumi{\arabic{enumi}.}
\setcounter{enumi}{1}
\item
  离散位移群可以含有反射，也可以不含有。早在 1852 年，默比乌斯就实质性地确定了球面几何中由反射生成的有限位移群（这等价于 \(\mathbf{R}^3\) 中有限欧几里得位移群的同一问题）；他发现，除去循环群之外，这样的群都以一个球面三角形为基本域，其角具有形式 \(\pi / p, \pi / q, \pi / r\)，其中 \(p, q, r\) 是三个 \(>1\) 的整数，满足 \(\frac{1}{p} + \frac{1}{q} + \frac{1}{r} > 1\)（{[}223{]}, v. II, pp.~349-360 与 pp.~561-708）。他还注意到，这些群作为子群包含着所有的有限位移群。
\item
  当继黎曼和施瓦茨关于超几何函数与共形表示的工作之后，开始研究用以为圆弧所界的图形对复平面或半平面所作的“铺砌”时，上述最后这股思想潮流获得了新的扩展；克莱因和庞加莱把它奠立为“自守函数”理论的基础，并在其中（就圆弧与一条固定直线正交的情形）认识到一个与求非欧双曲平面（等同于“庞加莱半平面”）位移群的离散子群等价的问题{[}118{]}。
\item
  正多面体的概念以及用这种多面体对 \(R^{3}\) 作铺砌的概念，被施莱夫利推广到了所有的欧几里得空间 \(R^{n}\)，这项工作可追溯到 1850 年前后，但很久以后才得以发表，并且在很长时间里一直无人问津（{[}273{]}, v. I, pp.~167-387）；他完全确定了每个 \(R^{n}\) 中的正“多胞形”、保持这种多胞形不变的位移群，以及该群的一个基本域，与默比乌斯研究过的 n = 3 的情形一样，这个基本域是一个“房间”，它在球面 \(S_{n-1}\) 上的迹是一个球面单形。然而，他没有着手处理求 \(R^{n}\) 中由反射生成的有限位移群这一逆问题；这个问题很久以后才得到解决：n = 4 的情形由古尔萨{[}132{]}解决，而对任意的 n，其解决则要等到 E. 嘉当（{[}52 a{]}, v. I₂, pp.~1003-1020）与考克斯特{[}70 c{]}的工作，我们稍后将再回到这些工作。
\end{enumerate}

1890 年前后，随着基灵和 E. 嘉当关于李群最初的工作，一股新的思想潮流开始了，它在很长时间里将与前面的那些工作互无联系地发展。基灵{[}180{]}和嘉当（{[}52 a{]}, v. I \(_{1}\), pp.~137-287）在研究复半单李代数的结构时，从一开始就打算让某些线性形式 \(\omega_{\alpha}\) 起着首要的作用，这些线性形式定义在这样一个李代数的一个“嘉当子代数”\(\mathfrak{h}\) 上，而这个李代数就是 \(\mathfrak{g}\)；它们就是相对于 \(\mathfrak{h}\) 的“根”，之所以如此称呼，是因为在基灵那里它们表现为特征方程 \(\det(\mathrm{ad}_{\mathfrak{g}}(x) - T) = 0\) 的根，这里把它们看作 \(x \in \mathfrak{h}\) 的函数。基灵和嘉当所建立的这些“根”的性质，用今天的几何语言来说，无非是断言它们构成一个“约化根系”；

他们随后证明，复半单李代数的分类可以归结为相伴“根系”的分类，而后者本身又可以归结为某些整系数矩阵（后来称为“嘉当矩阵”）的确定。基灵和嘉当还揭示出，对每个根 \(\omega_{\alpha}\)，存在根集合的一个对合置换 \(S_{\alpha}\)；他们本质性地使用了变换 \(C = S_{\alpha_{1}}S_{\alpha_{2}} \ldots S_{\alpha_{l}}\)，即与构成一个基础根系的 l 个根相伴的诸置换的乘积（这一变换今天称为“考克斯特变换”）；他们甚至把这个置换扩张为由基础根 \(\omega_{\alpha_{i}}(1 \leq i \leq l)\) 生成的向量空间上的一个线性变换，并研究其特征值（{[}180{]}, II, p.~20；{[}52 a{]}, v. I \(_{1}\), p.~58）。但无论是基灵，还是起初的嘉当，似乎都没有想到考虑群 \(G'\)，即由诸 \(S_{\alpha}\) 生成的群；而当嘉当稍后（{[}52 a{]}, v. I \(_{1}\), pp.~293-353）着手确定特征方程

\[
\det (\mathrm{ad} _ {\mathfrak {g}} (x) - T) = 0
\]

（即“一般元素”\(x \in h\) 的特征方程）的伽罗瓦群 G 时，他起初并没有援引 \(S_{\alpha}\) 来研究它；三十年后，已经处于 H. 外尔工作的影响之下，他证明了（{[}52 a{]}, v. I \(_{1}\), pp.~555-568）G 以群 \(G'\) 为正规子群，并在所有情形确定了商群 \(G/G'\) 的结构：对单代数 g 而言，它是 1 阶或 2 阶的，但 \(D_{4}\) 型除外，那时它同构于 \(S_{3}\)；也是在这次机会中，他把 \(G'\) 解释为复半单李代数中使一个嘉当子代数固定不动的内自同构所诱导的群。

我们刚才提到的 H. 外尔的工作，开创了群 \(G'\)（后来称为 g 的“外尔群”）的几何解释；正如基灵和嘉当对变换 C 所做的那样，他想到把 \(S_{\alpha}\) 看作 h 上线性形式向量空间中的反射。也是在 H. 外尔的论文{[}331 b{]}中，可以看到“仿射外尔群”的基本域的出现（不过它与“外尔群”\(G'\) 的联系并没有被明确指出）；外尔利用它证明了紧半单群的基本群是有限的，这是他证明复半单李代数线性表示完全可约性的一个关键点。不久之后，E. 嘉当完成了综合，把 H. 外尔的整体观点、他自己的实或复半单李代数理论以及他当时正在建立的对称黎曼空间理论融为一体。在论文（{[}52 a{]}, v. I₂, pp.~793-840）中，他完成了外尔群与仿射外尔群的基本多胞形的确定，并引入了权系和根基权系；他把这一讨论推广到对称空间，从而特别遇到了非约化根系的第一批例子（{[}52 a{]}, v. I₂, pp.~1003-1010）。最后，论文（{[}52 a{]}, v. I₂, pp.~867-989）给出了下列事实的第一个证明：\(R^{n}\) 中每个由反射生成的不可约有限群都具有一个基本域，它在球面 \(S_{n-1}\) 上的迹是一个球面单形；也正是在这项工作中，他借助几何的考虑证明了最大根的唯一性（在根系上任取的一个字典序下）。

稍后，范德瓦尔登{[}317 b{]}依靠 H. 外尔的论文证明了复半单李代数的分类等价于约化根系的分类，并通过初等的几何考虑完成了这一分类（而在基灵和嘉当那里，这一分类是复杂的行列式计算的结果）。大约同一时期，考克斯特明确地确定了所有由反射生成的不可约有限欧几里得位移群{[}70 c{]}；他就这样补全了 E. 嘉当论文（{[}52 a{]}, v. I₂, pp.~793-840）的结果，因为嘉当当时只确定了“晶体”群（即与一个根系相伴的群，或者也可说能够嵌入一个无限离散位移群的群）。第二年{[}70 d{]}，考克斯特证明了：由反射生成的有限群是仅有的（在相差一个同构的意义下）具有如下表出的有限群——由对合生成元 \(R_i\) 加上形如 \((R_i R_j)^{m_{ij}} = 1\) 的关系（\(m_{ij}\) 为整数）给出；从那时起，凡是具有这种表出的群（有限的或否）都得到了“考克斯特群”这一名称。

我们在上面描述过的两股研究潮流之间的第一个联系，看来是由考克斯特{[}70 e{]}、随后由维特{[}337 c{]}建立的。他们注意到，由反射生成的不可约无限欧几里得位移群与复单李代数（在相差一个同构的意义下）一一对应。维特给出了这类离散群的一个新的确定，并且还推广了上面回顾的考克斯特定理{[}70 d{]}，同时刻画了同构于无限离散欧几里得位移群的那些考克斯特群。这一结果，加上双曲几何中的类似群也是考克斯特群这一事实，\(^{1}\) 使人们得以公开着手研究这些群，先是把重点放在几何实现上（{[}71{]}、{[}308 a{]}），随后按照 J. 蒂茨{[}308 b 和 c{]}的做法在纯代数的框架内进行。

从维特的工作开始，半单李群的理论与由反射生成的离散群的理论将不断地彼此发生作用，而且成果极其丰硕。早在 1941 年，施蒂费尔{[}296{]}就指出，外尔群恰好是使一个格保持不变的有限反射群。谢瓦莱{[}62 e{]}和哈里希-钱德拉{[}149 a{]}在 1948 至 1951 年间对“晶体”群与复半单李代数之间的双射对应给出了先验的证明；在那之前，人们除了对每一类型的单李代数分别验证这一对应之外，别无他法。

另一方面，大约在 1950 年，人们注意到在外尔群下不变的多项式在无穷维线性表示的理论{[}149 a{]}以及李群的拓扑学中起着重要的作用。

就考克斯特而言，他{[}70 g{]}重新拾起了对变换 C（由反射生成的有限群 W 的基础反射的乘积）的研究，并注意到（通过对每个类型分别考察）W 下不变的多项式代数由代数无关的元素生成，这些元素的次数与 C 的特征值以简单的方式相联系。这些结果的先验证明随后由谢瓦莱{[}62 f{]}就第一个结果、由科尔曼{[}68{]}和斯坦伯格{[}292{]}就第二个结果给出。

随着 A. 博雷尔关于线性代数群{[}30{]}的工作，李群理论中开始了一些新的发展，这些发展后来导致了这一理论的显著扩展。A. 博雷尔揭示出李群中极大连通可解子群（后来称为“博雷尔子群”）的重要性，并把它们当作主要工具，用来把经典理论的一大部分移植到代数闭域上的代数群（不过当时还没有得到单代数群的分类\(^{2}\)）。（在实或复经典群的情形）博雷尔子群几年之前就已经出现在盖尔范德和纽马克关于无穷维表示的工作中；1954 年，F. 布吕阿发现了一个引人注目的事实：对经典单群而言，群按博雷尔群所作的双陪集分解以典范的方式由外尔群来标号{[}40{]}。这一结果随后由哈里希-钱德拉{[}149 b{]}推广到所有的实和复半单群。另一方面，1955 年，谢瓦莱{[}62 g{]}成功地做到了：对每个复半单李代数 g 和每个交换域 k，联系一个系数在 k 中的矩阵群，它具有布吕阿分解；他利用后一事实证明，除去少数例外，这样定义的群是单群（在抽象群理论的意义下）。他就这样“解释”了最迟自若尔当和李以来就已注意到的重合，即 A、B、C、D 型的单李群（在李群理论的意义下）与在任意域上以纯代数方式定义的经典单群之间的重合（直到那时，这一重合只能被推广到例外型 \(G_{2}\) 和 \(E_{6}\)，那是迪克森做到的{[}88c 和 d{]}）。特别地，取 k 为有限域时，谢瓦莱的构造对每一类型的复单李代数都给出一族有限单群，其中包含当时已知的有限单群的很大一部分，还包含三个新的系列（对应于单李代数 \(F_{4}, E_{7}\) 和 \(E_{8}\) 的类型）。不久之后，多位作者（赫齐格、铃木、Ree、斯坦伯格和蒂茨）用各种程序——即对谢瓦莱方法的种种修改——一方面证明了可以用类似的方式得到当时已知的其他有限单群（交错群与马蒂厄群除外），另一方面构造出了新的有限单群的其他系列（参见{[}54{]}）。

几乎在同一时期，谢瓦莱{[}62 h{]}仍然使用布吕阿分解的技术，再加上关于博雷尔子群正规化子的一个关键结果，重新研究了线性代数群，并得到了这样的结果：在任意特征的特征代数闭域 k 上，半单线性代数群\(^{3}\)的理论本质上给出与 k = C 时的基灵—嘉当分类相同的类型。接着，J. 蒂茨{[}308 a 和 b{]}通过分析谢瓦莱的方法，得到了布吕阿分解的一个公理化版本（“BN 对”），其形式异常灵活，只涉及群的结构；目前正是这一概念以“蒂茨系”的名称为人所知。我们在上面谈到的所有单群（该词各种意义下的单群）都典范地带有蒂茨系，而蒂茨本人{[}308 c{]}证明了：在抽象群 G 中这样一个系的存在，加上来自纯群论的几个补充性质，就使人们能够证明 G 是单群，这条定理涵盖了迄今为止对这些群所给出的绝大多数单性证明。另一方面，他与 A. 博雷尔合作，推广了{[}62 h{]}中谢瓦莱的结果，证明了在任意域上的半单线性代数群的有理点群中存在蒂茨系{[}31{]}。

这些问题中出现的所有蒂茨系都具有有限的外尔群。另一类例子由岩堀和松本{[}170{]}发现；他们证明了，如果在{[}62 g{]}中谢瓦莱的构造里取 k 为一个 p 进域，那么所得到的群带有一个蒂茨系，其外尔群是出发时所用的那个复半单李代数的仿射外尔群。这一结果刚刚由布吕阿和蒂茨{[}41{]}推广到局部域上的所有半单代数群。

